\uuid{FlPS}
\exo7id{685}
\titre{exo7 685}
\auteur{gourio}
\organisation{exo7}
\datecreate{2001-09-01}
\isIndication{false}
\isCorrection{false}
\chapitre{Continuité, limite et étude de fonctions réelles}
\sousChapitre{Continuité : pratique}
\cpge{mpsi-mp2i.fonctions-valeurs-reelles-complexes.a, pcsi.techniques-calcul-fonctions.a, ptsi.techniques-calcul-fonctions.a, mpsi-mp2i.limites-continuite.b, pcsi.limites-continuite.b, ptsi.limites-continuite.b, mpsi-mp2i.limites-continuite.c, pcsi.limites-continuite.c, ptsi.limites-continuite.c, bcpst1.fonctions-reelles-usuelles.a, bcpst1.limites-continuite.c, pc.espaces-vectoriels-normes.f, psi.espaces-vectoriels-normes.f}
\module{Analyse}
\niveau{L1}
\difficulte{}

\contenu{
\question{
Soit $f$ d\'{e}finie par $f(x)=E(x)+\sqrt{x-E(x)}, $ o\`{u} $E$ d\'{e}signe
la partie enti\`{e}re.
Donner le domaine de d\'{e}finition de $f,$ puis une relation entre $f(x+1)$
et $f(x)$. $f$ est-elle monotone?\ $f$ est-elle $k-$lipschitzienne sur $[a,1] $($a>0$) ?
Et sur $[0,1]$ ?
\'Etudier la continuit\'{e} de $f$ sur $[0,1]$ en utilisant la
d\'{e}finition. D\'{e}duisez en la continuit\'{e} sur $\Rr$.
}
}
